\chapter{Chest Tube Insertion}

\section*{Overview}
Chest tube insertion, or tube thoracostomy, is placement of a drain into the pleural space to evacuate air, blood, pus, or fluid and to re-expand the lung. The exact approach, setting, and preparation depend on the indication, anatomy, and the patient's health.

\section*{Alternative Names}
Tube thoracostomy; intercostal drain placement; pleural drain

\section*{Principle}
A drainage catheter is inserted into the pleural space through the chest wall and connected to a sealed drainage system with an underwater seal. The seal allows one-way outflow of air or fluid while the lung expands.

\section*{History}
Chest drainage has been practiced for empyema since the 19th century, popularized by surgeons during World War I. Closed underwater-seal drainage became standard during the Korean and Vietnam wars.

\section*{Why It Is Done}
Performed for pneumothorax (especially tension or large spontaneous), hemothorax, pleural effusion or empyema, post-thoracotomy drainage, and after cardiothoracic surgery.

\section*{Is This a Primary or Secondary Test or Procedure}
Primary emergency and therapeutic procedure for pleural air and fluid collections.

\section*{How It Is Performed}
Performed at the bedside or in the operating room. The patient is positioned, the skin over the fourth or fifth intercostal space in the midaxillary line is prepped and anesthetized, a small incision is made, the tract is bluntly dissected, and the tube is guided into the pleural space toward the apex (for air) or base (for fluid). It is secured and connected to suction or an underwater seal.

\section*{Preparation}
Consent, coagulation assessment, and in elective cases imaging to localize the collection. Emergency insertion may proceed without delay.

\section*{Contraindications and Precautions}
Relative contraindications include coagulopathy, bleeding diathesis, and local infection at the site. Precaution: do not insert above the level of the nipple line without imaging guidance to avoid injury to the subclavian vessels.

\section*{Risks}
Injury to lung, intercostal vessels, liver, spleen, or diaphragm; bleeding; infection; malposition; and re-expansion pulmonary edema.

\section*{Aftercare and Recovery}
The tube is connected to a drainage system, serial chest radiographs confirm position and lung expansion, and the tube is removed when the air leak or drainage resolves and the lung remains expanded.

\section*{Outcome and Follow-up}
Chest radiograph confirms tube position and lung re-expansion. Output volume and air leak guide removal timing; laboratory analysis of drained fluid may guide diagnosis.

\section*{Expected Results}
Complete lung re-expansion, cessation of air leak, and drainage of the collection, with tube removal after 24-48 hours of stability.

\section*{Follow-up}
Chest radiograph after removal and clinical follow-up for recurrence of pneumothorax or effusion.

\section*{When to Seek Medical Care}
Follow the procedure team's specific instructions. Seek urgent medical assessment for severe or worsening pain, uncontrolled bleeding, fever or signs of infection, trouble breathing, chest pain, fainting, new weakness or confusion, inability to urinate, or any rapidly worsening symptom. The appropriate warning signs depend on the procedure and the patient's medical condition.

\section*{References}
\begin{enumerate}
  \item MedlinePlus. Chest tube insertion. U.S. National Library of Medicine; 2025.
  \item Current Medical Diagnosis and Treatment. McGraw-Hill; 2025.
\end{enumerate}

\clearpage
